\section{Introduction}

\subsection{The Integration Challenge}

AI applications increasingly depend on heterogeneous capabilities provided by model vendors, external services, and local devices. Providers such as OpenAI~\cite{openai_docs}, Anthropic~\cite{anthropic_docs}, Google, and DeepSeek expose powerful models through APIs, but direct integration introduces tight coupling and escalating system complexity.

Most current AI frameworks address this challenge through application-level orchestration (e.g., LangChain~\cite{langchain_docs}). Although this approach improves development velocity, it frequently combines orchestration policy and execution mechanics in one runtime, thereby increasing lock-in and reducing component reusability.

\subsection{The Missing Layer}

The ecosystem still lacks a minimal execution layer dedicated to capability invocation. Such a layer should:

\begin{enumerate}[leftmargin=1.5em]
    \item Execute capabilities defined by standardized protocols.
    \item Translate protocol requests into provider-specific API calls.
    \item Normalize results into standardized responses.
    \item Handle execution errors without implementing application-level retry strategies.
\end{enumerate}

This layer excludes higher-level concerns such as planning, workflow orchestration, and application logic, thereby preserving a lightweight and reusable execution boundary.

\subsection{Contributions}

The primary contributions are as follows:

\begin{enumerate}[leftmargin=1.5em]
    \item \textbf{Formal Architectural Model:} A three-layer logical model $S = (A, P, E)$ is specified to separate cognition, orchestration, and execution with explicit layer responsibilities.
    \item \textbf{Minimal Runtime Constraint:} Minimality is formalized as a verifiable design constraint, with explicit runtime responsibilities and non-responsibilities for deterministic capability execution.
    \item \textbf{Unified Capability Semantics:} Capability is formalized as $C = (I, O, \Sigma, Q, \Pi)$, enabling technical invocation and economic observability to share a single schema-level object.
    \item \textbf{Cross-Layer Integration:} The three-layer logical model is mapped to the five-layer AI-Native implementation stack, clarifying how execution-layer invariants are preserved in full-system deployments.
\end{enumerate}

The execution layer is then positioned within the broader AI-Native architecture developed in later sections. Figure placeholders are introduced in this initial conversion pass; the ASCII diagrams from the Markdown source will be redrawn as vector figures in the next phase.
